% §VI --- Experiments. Status: REWRITTEN on the frozen v8 environment.
% All numbers from outputs/rl/eps_recal_v8 + flat_ppo_v8 + m3_ablation_v8 +
% m4_generalization_v8 + separation_v2 (paper2_numbers_v8.json).
\section{Experiments}\label{sec:experiments}

We follow the reporting template of constrained-DRL
practice~\cite{khairy2021constrained,long2025lyahippo}: a convergence
triptych with the dual trajectories, a dual-condition main comparison with
oracle upper bounds and non-learning baselines, per-channel constraint
accounting, a dual-machinery ablation, four-axis generalization sweeps
with a dedicated zero-shot section, and the corrected safety bridge.

\subsection{Protocol and metrics}\label{sub:protocol}
All runs use the frozen environment of
Sections~\ref{sec:system}--\ref{sec:algorithm} (communication-window
deadlines, attainability-calibrated $\varepsilon_k$, certified escalation,
structural cache ban, mission-aligned reward attribution, fair-share
reference bandwidth, LUT out-of-support outage guard, per-channel dual
caps with warm-up). Learning arms train $500$ episodes $\times3$ seeds and
are evaluated deterministically on $50$ held-out episodes ($1{,}000$
queries) per condition; non-learning baselines use the same evaluation.
The peak-condition escalation budget is
$\delta_{\mathrm{esc}}=0.155$ and the nominal budget $0.05$
(Section~\ref{subsec:cert}). We report: \textbf{mission} (quality- and
deadline-compliant fraction of the \emph{admitted} set, the honest
headline of Eq.~\eqref{eq:cmdp}), \textbf{task} (end-to-end task success),
\textbf{acc} (mean answer accuracy on the admitted set), \textbf{cacheR}
(cache-route share), \textbf{neRejR} (non-escalated rejects: serveable
queries the policy dropped), \textbf{critEsc} (critical-escalation rate
vs.\ its budget), \textbf{ddlVio} (admitted-set deadline violations), and
\textbf{conflict} (tactical-conflict rate vs.\ the $0.08$ TLS budget).

Two oracles bound the problem from above. The \emph{best-feasible oracle}
picks, per query, the cheapest evidence level that succeeds end-to-end,
but never escalates; the \emph{escalation-aware oracle} additionally
honours the certificate---a critical query with no spec-attainable
transmission service is escalated rather than served badly. The gap
between them measures what admission structure is worth.

\subsection{Main dual-condition comparison}\label{sub:main}
Table~\ref{tab:main} reports the main comparison.

\textbf{Peak is an admission regime.} The escalation-aware oracle reaches
admitted mission $0.916$ (admitted deadline violations $0.084$)---but only
by escalating $0.666$ of the critical load. The best-feasible oracle,
denied escalation, collapses to $0.528$ with admitted deadline violations
of $0.412$: \emph{even an oracle cannot keep the mission constraint on the
admitted set without shedding load}. The feasible region for a
non-escalating strategy is essentially empty at peak, and safety in this
regime is created by the certificate-plus-escalation structure, not by any
strategy-layer optimization. Among learning arms, the constrained and
fixed-penalty policies land at admitted mission
$0.180$/$0.184$ with near-zero conflict ($0.003$, $27\times$ under
budget) and tight admitted deadline compliance ($0.023$), but they
over-reject (neRejR $\approx0.35$)---dropping serveable queries is the
cheap escape from a pinned quality price
(Section~\ref{sub:pathology})---while the unconstrained arm escalates
$2$--$4\times$ more ($0.329$, over budget) and rejects less. The flat
single-head baseline is instructive in the opposite direction: it trains
stably ($0.344$ admitted mission) but \emph{never interacts with the
admission machinery at all} (zero rejects, zero escalations)---it serves
structurally unserveable critical queries rather than handing them up,
its admitted quality-violation rate is $0.564$, and its conflict rate is
$22\times$ the proposed policy's ($0.065$ vs.\ $0.003$, still under
budget). No learning arm approaches the escalation-aware ceiling; we
return to what peak \emph{does} certify in
Section~\ref{sub:channels}.

\textbf{Nominal is a pricing regime.} Under the feasible nominal
condition the proposed policy reaches admitted mission $0.653$---$93\%$ of
the escalation-aware oracle's $0.703$ and above semantic-greedy's
$0.626$---with zero escalations and zero conflicts; trained directly on
nominal it reaches $0.723$, and the fixed-penalty arm (its static quality
price clamped to the same per-channel cap) reaches $0.727$. The flat
baseline trails the hybrid arms ($0.592$), and the unconstrained arm gives
up deadline discipline (admitted deadline violations $0.254$ vs.\ the
proposed $0.168$ and nominal-trained $0.047$). The backend ablation
(legacy LUT quality instead of the calibrated v5 table) shifts absolute
numbers ($0.443$ nominal) without changing any ordering, so claims are
conditional on the calibration but not artifacts of it.

\begin{table*}[t]
\centering
\caption{Main dual-condition comparison on the frozen v8 environment
(mean$\pm$std over $3$ seeds where applicable; oracles and heuristics are
deterministic). Mission/acc/ddlVio are on the admitted set
(Eq.~\eqref{eq:cmdp}); critEsc is against the escalation budget
($0.155$ peak / $0.05$ nominal); conflict is against the $0.08$ TLS
budget.}
\label{tab:main}
\scriptsize
\setlength{\tabcolsep}{2.2pt}
\begin{tabular}{@{}llcccccccc@{}}
\toprule
Cond. & Method & Mission (adm.) & Task & Acc (adm.) & CacheR & neRejR & CritEsc & DdlVio (adm.) & Conflict \\
\midrule
\multirow{9}{*}{peak}
 & Esc-aware oracle & $0.916$ & $0.789$ & $0.808$ & $0.000$ & $0.000$ & $0.666$ & $0.084$ & $0.193$ \\
 & Best-feasible oracle & $0.528$ & $0.260$ & $0.757$ & $0.008$ & $0.000$ & $0.000$ & $0.412$ & $0.612$ \\
 & Semantic greedy & $0.616$ & $0.268$ & $0.719$ & $0.008$ & $0.000$ & $0.000$ & $0.380$ & $0.612$ \\
 & Always cache & $0.000$ & $0.000$ & $0.000$ & $1.000$ & $0.000$ & $0.000$ & $0.000$ & $0.000$ \\
 & Random & $0.292\pm0.007$ & $0.120\pm0.017$ & $0.454\pm0.010$ & $0.339\pm0.019$ & $0.000$ & $0.000$ & $0.333\pm0.022$ & $0.419\pm0.027$ \\
 & Flat PPO (single head) & $0.344\pm0.051$ & $0.295\pm0.052$ & $0.361\pm0.021$ & $0.549\pm0.026$ & $0.000$ & $0.000$ & $0.093\pm0.036$ & $0.065\pm0.023$ \\
 & No-Lagrangian & $0.164\pm0.015$ & $0.384\pm0.226$ & $0.213\pm0.074$ & $0.298\pm0.193$ & $0.287\pm0.076$ & $0.329\pm0.271$ & $0.093\pm0.097$ & $0.005\pm0.002$ \\
 & Fixed-penalty & $0.184\pm0.013$ & $0.214\pm0.024$ & $0.168\pm0.007$ & $0.452\pm0.029$ & $0.327\pm0.040$ & $0.081\pm0.016$ & $0.021\pm0.009$ & $0.003\pm0.002$ \\
 & \textbf{Proposed} & $0.180\pm0.011$ & $0.210\pm0.023$ & $0.166\pm0.005$ & $0.439\pm0.035$ & $0.345\pm0.042$ & $0.086\pm0.021$ & $0.023\pm0.010$ & $0.003\pm0.002$ \\
\midrule
\multirow{9}{*}{nominal}
 & Esc-aware oracle & $0.703$ & $0.703$ & $0.454$ & $0.000$ & $0.033$ & $0.000$ & $0.234$ & $0.000$ \\
 & Best-feasible oracle & $0.691$ & $0.691$ & $0.485$ & $0.037$ & $0.000$ & $0.000$ & $0.281$ & $0.000$ \\
 & Semantic greedy & $0.626$ & $0.626$ & $0.461$ & $0.052$ & $0.000$ & $0.000$ & $0.339$ & $0.000$ \\
 & Always cache & $0.250$ & $0.250$ & $0.000$ & $0.750$ & $0.000$ & $0.000$ & $0.000$ & $0.000$ \\
 & Flat PPO (single head) & $0.592\pm0.074$ & $0.592\pm0.074$ & $0.441\pm0.018$ & $0.162\pm0.054$ & $0.000$ & $0.000$ & $0.263\pm0.112$ & $0.000$ \\
 & No-Lagrangian & $0.593\pm0.072$ & $0.593\pm0.072$ & $0.438\pm0.048$ & $0.151\pm0.083$ & $0.000$ & $0.000$ & $0.254\pm0.148$ & $0.000$ \\
 & Fixed-penalty & $0.727\pm0.017$ & $0.727\pm0.017$ & $0.409\pm0.011$ & $0.155\pm0.017$ & $0.000$ & $0.000$ & $0.049\pm0.031$ & $0.000$ \\
 & \textbf{Proposed} & $0.653\pm0.097$ & $0.653\pm0.097$ & $0.436\pm0.022$ & $0.153\pm0.067$ & $0.000$ & $0.000$ & $0.168\pm0.159$ & $0.000$ \\
 & Proposed (nominal-trained) & $0.723\pm0.017$ & $0.723\pm0.017$ & $0.421\pm0.010$ & $0.170\pm0.060$ & $0.000$ & $0.000$ & $0.047\pm0.028$ & $0.000$ \\
\bottomrule
\end{tabular}
\end{table*}

\subsection{Per-channel constraint accounting}\label{sub:channels}
Table~\ref{tab:channels} reports every dual channel's realized evaluation
cost against its limit, in the implemented constraint semantics
(joint risk-class event rate on the admitted set). The pattern is the
paper's core empirical claim. \emph{Structural channels hold everywhere}:
conflict ($0.003$ peak, $0.000$ nominal, vs.\ $0.08$), battery, GPU, both
critical-deadline channels at peak, and the escalation budget ($0.086$ vs.\
$0.155$ peak; $0$ vs.\ $0.05$ nominal). \emph{Quality channels expose the
regime split}: at peak the realized quality cost sits an order of
magnitude above its limit ($0.41$ vs.\ $0.02$ critical)---the infeasible
domain, where the accounting shows exactly \emph{which} promise cannot be
made and the certificate/escalation layer bounds what happens to the
unserveable load instead. Under nominal the quality costs drop to
$0.08$--$0.12$: still above the tight risk-aware limits at episode $500$,
with an interior, still-climbing dual pricing the residual gap
(Section~\ref{sub:pathology})---approaching, not attained, and stated as
such. We emphasize what this table replaces: an earlier draft of this work
claimed ``constraint satisfaction by construction''; the honest,
channel-resolved statement is that the \emph{hard-structure} channels are
satisfied by construction, the escalation budget is satisfied by
certificate, and the quality channels are priced but not closed.

\begin{table}[t]
\centering
\caption{Per-channel realized eval cost vs.\ limit $d_i$ (admitted set;
$\checkmark$ within limit, $\times$ above). Peak/nominal columns: the
peak-trained policy in its own condition and evaluated on nominal;
nom-tr = trained on nominal.}
\label{tab:channels}
\footnotesize
\setlength{\tabcolsep}{4pt}
\begin{tabular}{@{}lcccc@{}}
\toprule
Channel $c_i$ & $d_i$ & peak & nominal & nom-tr \\
\midrule
quality (normal)   & $0.05$ & $0.389\,\times$ & $0.101\,\times$ & $0.109\,\times$ \\
quality (critical) & $0.02$ & $0.408\,\times$ & $0.083\,\times$ & $0.121\,\times$ \\
deadline (normal)  & $0.05$ & $0.022\,\checkmark$ & $0.090\,\times$ & $0.001\,\checkmark$ \\
deadline (critical)& $0.02$ & $0.001\,\checkmark$ & $0.077\,\times$ & $0.046\,\times$ \\
conflict           & $0.08$ & $0.003\,\checkmark$ & $0.000\,\checkmark$ & $0.000\,\checkmark$ \\
battery            & ---    & $0.000$ & $0.000$ & $0.000$ \\
GPU (VRAM)         & ---    & $0.000$ & $0.000$ & $0.000$ \\
escalation         & $0.155/0.05$ & $0.086\,\checkmark$ & $0.000\,\checkmark$ & $0.000\,\checkmark$ \\
\bottomrule
\end{tabular}
\end{table}

\subsection{Convergence and the dual-pathology contrast}\label{sub:pathology}
Fig.~\ref{fig:conv} shows the convergence triptych---episode return, the
quality-critical constraint cost against its $0.02$ limit, and the dual
trajectories---for the peak-trained and nominal-trained proposed policy
($3$ seeds, $\pm$std). The dual warm-up (episodes $0$--$149$, shaded)
freezes ascent while the BC/prior shaping decays; what happens after it
lifts is the paper's central diagnostic, isolated in
Fig.~\ref{fig:lambda}. Under \emph{peak}, $\lambda_{QC}$ resumes from
$\approx\!0$ and ratchets monotonically---$1.95$ by episode $200$, $5.48$
by $300$, pinned at the per-channel cap $8.00$ by episode
$\approx\!370$ on $3/3$ seeds with zero fall-back---against a realized
cost stuck at $\approx\!0.4$. Under \emph{nominal}, the same machinery
climbs $\approx\!4\times$ slower and stays interior: $0.50$ at episode
$200$, $1.33$ at $300$, $3.29$ at $500$, never touching the cap, against
a cost declining toward its limit ($0.053$ at the training tail). One
algorithm, one schedule, two textbook behaviours: an interior shadow
price where the constraint is feasible, and the classical divergence of
the multiplier where it is not~\cite{altman1999constrained,chow2017risk}.
The pin is robust across three engineering generations (caps $20/20/8$
pin at $18.3/18.4/8.0$) and across the quality backend ($2/2$ seeds under
the legacy LUT), while the constrained arm's peak behaviour is invariant
under all of it (admitted mission $0.179/0.179/0.180$)---the levers move
the ceiling, not the feasibility.

\evfigwide{fig1_convergence.pdf}{0.98}{Convergence triptych on the frozen v8
environment: (a) episode return, (b) quality-critical constraint cost vs.\
its $0.02$ limit, (c) dual trajectories ($\lambda_{QC}$, and
$\lambda_{\mathrm{conflict}}$ at peak), peak-trained vs.\ nominal-trained,
$3$ seeds $\pm$std; dual warm-up window shaded.}{fig:conv}

\evfig{fig2_lambda_golden.pdf}{0.85}{The dual-pathology contrast:
$\lambda_{QC}$ ratchets to its cap under peak (infeasibility signature)
and converges interior under nominal (shadow price), same machinery and
schedule.}{fig:lambda}

\subsection{Dual-machinery ablation}\label{sub:m3}
Table~\ref{tab:ablation} ablates the dual machinery itself at peak
($3$ seeds each). \textbf{(i) Quality channel off}
($\lambda^{\max}_{Q}{=}0$): removing the quality price entirely leaves the
admitted-set mission statistically unchanged ($0.155$ vs.\ $0.180$)---in
the infeasible domain the price was not buying compliance---but the
behavioural split moves toward the unconstrained arm (cache $0.304$,
escalation $0.230$ and over budget): the pinned price had been shaping
\emph{how} the policy sheds load, not \emph{whether} it satisfies quality.
\textbf{(ii) Warm-start at the pin} ($\lambda_{QC}(0){=}8$, no warm-up):
starting the dual where it will end leaves the headline unchanged
(admitted mission $0.177$ vs.\ $0.180$; $\lambda_{QC}$ enters at $8.00$
and never leaves on $3/3$ seeds)---the pin is the domain's destiny, not an
artifact of the ascent schedule. The one visible effect is again on the
load-shedding split: with the full price active from episode $0$ the
critical-escalation rate lands at $0.134$, \emph{inside} the
certificate band $[0.105,0.205]$, where the warmed-up reference
under-escalates ($0.086$)---corroborating that in an infeasible domain the
dual's schedule shapes \emph{how} load is shed, not whether the constraint
is met.
\textbf{(iii) No slow head}: removing the learned mobility actor
(deterministic mobility defaults) also leaves the headline unchanged
($0.176$) and even zeroes the residual conflict rate ($0.000$ vs.\
$0.003$; both $\ll0.08$)---in this admission-dominated regime the mobility
hierarchy is \emph{not} what suppresses conflict, and we retract the
stronger causal reading an earlier draft attached to it. Its visible
effect is, once more, on the shedding split (critical escalation $0.168$,
the top of the band). The three ablations tell one story: at peak,
\emph{no} strategy-layer knob moves the admitted-mission headline (all
arms sit in $0.155$--$0.184$); each knob only redistributes how the
unserveable load is shed among escalate/reject/cache. The controlling
structure is the certificate and its budget---which is exactly what the
admission-regime reading predicts.

\begin{table*}[t]
\centering
\caption{Dual-machinery ablation at peak (mean$\pm$std over $3$ seeds).
$\lambda_{QC}^{\mathrm{end}}$ lists the terminal quality-critical dual per
seed; the reference row is the proposed v8 arm.}
\label{tab:ablation}
\scriptsize
\setlength{\tabcolsep}{3pt}
\begin{tabular}{@{}lcccccccc@{}}
\toprule
Arm & Mission (adm.) & Task & CacheR & neRejR & CritEsc & DdlVio (adm.) & Conflict & $\lambda_{QC}^{\mathrm{end}}$ \\
\midrule
Proposed (reference) & $0.180\pm0.011$ & $0.210\pm0.023$ & $0.439\pm0.035$ & $0.345\pm0.042$ & $0.086\pm0.021$ & $0.023\pm0.010$ & $0.003\pm0.002$ & $8.00/8.00/8.00$ \\
\;(i) quality channel off & $0.155\pm0.025$ & $0.276\pm0.115$ & $0.304\pm0.144$ & $0.392\pm0.044$ & $0.230\pm0.213$ & $0.051\pm0.036$ & $0.004\pm0.004$ & $0.00/0.00/0.00$ \\
\;(ii) warm-start at the pin & $0.177\pm0.016$ & $0.229\pm0.004$ & $0.444\pm0.012$ & $0.320\pm0.015$ & $0.134\pm0.037$ & $0.021\pm0.006$ & $0.003\pm0.002$ & $8.00/8.00/8.00$ \\
\;(iii) no slow mobility head & $0.176\pm0.017$ & $0.246\pm0.015$ & $0.435\pm0.025$ & $0.316\pm0.018$ & $0.168\pm0.057$ & $0.018\pm0.002$ & $0.000$ & $8.00/8.00/8.00$ \\
\bottomrule
\end{tabular}
\end{table*}

\subsection{Sample efficiency}\label{sub:sampleeff}
Table~\ref{tab:sampleeff} reports episodes to $90\%$ of the terminal
smoothed return (min-to-terminal span) on a uniform $500$-episode budget.
The BC warm start front-loads the peak arms (the return span is nearly
flat from episode $0$), so their ep-to-$90\%$ is small by construction; the
nominal-trained arm, whose return genuinely improves over training,
reaches $90\%$ at $146\pm123$ episodes, and the flat baseline at
$27\pm18$. Wall-clock is $\approx\!3$ minutes per $500$-episode arm on a
single consumer GPU---the whole $36$-arm matrix reruns in under an hour,
which is what made the three-generation lineage of
Section~\ref{sub:pathology} affordable.

\begin{table}[t]
\centering
\caption{Sample efficiency: episodes to $90\%$ of terminal smoothed
return; $500$-episode budget, $3$ seeds.}
\label{tab:sampleeff}
\footnotesize
\begin{tabular}{@{}lcc@{}}
\toprule
Arm & ep-to-$90\%$ & budget \\
\midrule
Proposed (peak) & $3\pm1$ & $500$ \\
No-Lagrangian (peak) & $2\pm1$ & $500$ \\
Fixed-penalty (peak) & $3\pm1$ & $500$ \\
Legacy-LUT backend (peak) & $2\pm2$ & $500$ \\
Flat PPO (peak) & $27\pm18$ & $500$ \\
Quality-channel-off (peak) & $3\pm2$ & $500$ \\
Warm-start-at-pin (peak) & $2\pm2$ & $500$ \\
No-slow-head (peak) & $3\pm2$ & $500$ \\
Proposed (nominal-trained) & $146\pm123$ & $500$ \\
\bottomrule
\end{tabular}
\end{table}

\subsection{Generalization and zero-shot robustness}\label{sub:m4}
Following the constrained-DRL reporting
standard~\cite{khairy2021constrained}, we evaluate the peak-trained
policies \emph{without retraining} along four test-time axes
(Fig.~\ref{fig:m4axes}; $5$ points each, $50$ rollouts per point,
$3$ seeds per learning arm) and at three unseen operating points
(Fig.~\ref{fig:zeroshot}).

\textbf{Fleet size ($2$--$8$ UAVs).} Admitted mission is flat for every
arm (proposed $0.172$--$0.191$; escalation-aware oracle
$0.909$--$0.916$), and the arm ordering is preserved at every point. The
padded observation and mask machinery transfer across fleet sizes without
retraining.

\textbf{Arrival load ($12$--$28$ queries per fixed burst horizon).} The
one axis with real motion, and it moves exactly as the admission-regime
reading predicts: the learning arms decline smoothly (proposed
$0.254\to0.134$) as the un-serveable fraction grows, while the
escalation-aware oracle's admitted mission \emph{rises} ($0.741\to0.933$)
because its escalation rate climbs with overload ($0.450\to0.667$)---the
harder the domain, the more the certificate sheds and the cleaner the
admitted core. The widening scissors between the oracle and the learning
arms is the clearest picture of what strategy-layer learning does not
solve at peak.

\textbf{Link attenuation ($+0$ to $+15.5$\,dB excess loss) and
interference floor ($-120\to-116$\,dBm).} Both link axes leave admitted
mission flat (proposed $0.180\to0.174$ and $0.180\to0.179$
respectively)---the extra attenuation drops the mean sensed SNR from
$28.8$ to $15.9$\,dB, still inside the calibrated quality plateau, and the
raised noise floor shifts SINR by $\lesssim\!2$\,dB. At the peak operating
point the bottleneck is the admission geometry, not the link: a useful
robustness statement, and an honest one---these axes would bite only past
the LUT support edge, where the outage guard (Section~\ref{sec:system})
takes over.

\textbf{Zero-shot triple.} At the \emph{unseen soft-conflict profile}
(calibrated softer UTM geometry) the peak-trained constrained arms
transfer well (proposed $0.619\pm0.219$, fixed-penalty $0.714\pm0.153$)
and beat the unconstrained arm ($0.473$)---and, notably, the
escalation-aware oracle itself ($0.486$), whose certificate-triggered
shedding is too aggressive for a near-feasible domain. At
\emph{$1.5\times$ load} the arms continue the arrival-axis trend
($0.123$). The \emph{unseen low-SNR blockage profile} is the failure
mode we highlight: the certificate marks essentially the whole critical
load unserveable---the escalation-aware oracle escalates nearly
everything and its admitted mission collapses to $0.000$---while the
zero-shot policies never engage escalation at all (critEsc $0.000$) and
keep serving with $0.22$ admitted deadline violations (proposed
$0.149\pm0.085$). The certificate transfers; the learned \emph{shedding
behaviour} does not. Making the escalate/serve split itself robust to
unseen link regimes is the concrete open problem this sweep isolates.

\evfigwide{fig3_m4_axes.pdf}{0.9}{Zero-shot generalization of the
peak-trained policies along four axes (no retraining): (a) fleet size,
(b) arrival load on the fixed burst horizon, (c) excess link attenuation,
(d) interference floor. Bands: $\pm$std over $3$ seeds.}{fig:m4axes}

\evfig{fig4_zeroshot.pdf}{1.0}{Three-point zero-shot evaluation: unseen
soft-conflict profile, $1.5\times$ arrival load, unseen low-SNR blockage
profile.}{fig:zeroshot}

\subsection{Safety bridge, corrected: evidence choice as an airspace tax}
\label{sub:bridge}
Finally we close the loop between communication and separation with the
M/G/1 coupling of Section~\ref{subsec:mg1}
(Table~\ref{tab:bridge}, Fig.~\ref{fig:bridge}). At the $-5$\,dB Rayleigh
worst case with the WP5 $2\sigma$ band, a token upload adds $7.5$\,ms of
queueing wait to the tactical loop: $\dtc=370.6$\,m, indistinguishable
from the BUBBLES default-comm baseline of $370.5$\,m, at $99.98\%$
relative capacity. An image upload at peak shared-C2 load adds $1.98$\,s:
$\dtc=394.3$\,m ($+23.7$\,m over baseline) and $94.0\%$ relative capacity
($\approx\!6\%$ of airspace forfeited); under nominal load the tax is
$+11.2$\,m and $\approx\!3\%$. A dedicated C2 band roughly $\!$halves the
peak image tax ($+17.2$\,m). \emph{Honest correction}: an earlier version
of this bridge substituted the payload airtime for the official $1.8$\,s
separation-communication mean, so light evidence appeared to
\emph{shorten} $\dtc$ below the deliverable's baseline ($-48.8$\,m,
``$1.14\times$ capacity''); that gain was an artifact of letting the
payload stand in for the separation loop's own exchange. Under the
corrected queueing coupling nothing beats the baseline---heavy evidence
taxes the airspace, light evidence avoids the tax---and the co-design
signal survives in honest form: at low SNR, routing to tokens keeps the
separation service at its certified baseline while the image path costs
up to $6\%$ of capacity precisely when spectrum is scarcest.

\begin{table}[t]
\centering
\caption{Corrected safety bridge at $-5$\,dB Rayleigh, $2\sigma$ band,
shared C2: M/G/1 wait $W$, contact distance, and relative capacity
(baseline $\dtc$: $370.49$\,m; dedicated-C2 peak image in parentheses).}
\label{tab:bridge}
\footnotesize
\setlength{\tabcolsep}{4pt}
\begin{tabular}{@{}lcccc@{}}
\toprule
Evidence / load & $W$ (s) & $\dtc$ (m) & $\Delta\dtc$ (m) & rel.\ capacity \\
\midrule
$s_1$ token, nominal & $0.004$ & $370.53$ & $+0.04$ & $1.000$ \\
$s_1$ token, peak    & $0.008$ & $370.58$ & $+0.09$ & $1.000$ \\
$s_2$ image, nominal & $0.930$ & $381.65$ & $+11.2$ & $0.971$ \\
$s_2$ image, peak    & $1.984$ & $394.29$ & $+23.8$ & $0.940$ \\
\;\;(dedicated C2)   & $1.435$ & $387.70$ & $+17.2$ & $0.956$ \\
\bottomrule
\end{tabular}
\end{table}

\evfigwide{fig8_safetybridge.pdf}{0.92}{Corrected safety bridge (Rayleigh,
$2\sigma$): (a) tactical-contact separation $\dtc$ and (b) relative
airspace capacity vs.\ SNR for token and image under peak/nominal shared-C2
load. The M/G/1 coupling adds queueing wait to the official $1.8$\,s
mean, so no evidence level beats the BUBBLES baseline; images pay an
airspace tax that grows as the link weakens.}{fig:bridge}
